\section{创业结构：三人平分股份、融资和第一笔大单}

上一章讲技术，本章看组织和资本。禾赛早期一个罕见设定是三人平分股份。对创业公司来说，这既是强信任，也带来治理要求。硬科技创业周期长、资本需求高、技术和商业不确定性大，如果创始人之间不能长期绑定，很难穿过早期低谷。

\begin{figure}[H]
\centering
\includegraphics[width=0.94\textwidth]{three-founder-equity.png}
\caption{三人平分股份：罕见股权结构背后是信任、能力互补和长期绑定。自制概念图，依据 00:32:13--00:43:35 对谈内容整理。}
\end{figure}

\begin{knowledgebox}{读图：平分股份不是浪漫主义}
三人平分可以强化共同体，但要求创始团队有足够强的互信、沟通和决策机制。否则一旦进入融资、转型和利益分配阶段，治理风险会放大。
\end{knowledgebox}

\subsection{融资不是故事，而是现金流桥梁}

前面讲股权结构，本节转向资本节奏，并回答硬科技公司为什么不能只会讲愿景。硬科技公司融资不只是讲愿景，它要支撑研发、样机、供应链、交付和客户验证。李一帆谈到融资节奏和现金流压力，说明硬科技创业不能只靠资本故事，最终要过订单和回款关。

\begin{figure}[H]
\centering
\includegraphics[width=0.96\textwidth]{financing-cashflow.png}
\caption{融资与现金流：硬件创业不能只讲故事，必须跨过样机、订单和现金流。自制概念图，依据 00:43:35--00:49:02 与 01:10:06--01:20:47 对谈内容整理。}
\end{figure}

\begin{knowledgebox}{读图：融资只是在验证链中买时间}
图中从技术故事到量产扩张，每一步都更接近真实市场。融资可以支撑研发，但订单和回款才说明客户愿意承担真实风险；量产则说明组织和供应链都开始过关。
\end{knowledgebox}

\subsection{第一笔 2000 万大单}

融资只能买时间，订单才能证明市场，本节看第一笔 2000 万大单的意义。第一笔 2000 万大单让团队第一次被市场验证，也把交付压力放大。对硬科技公司而言，大单不是单纯收入，而是信用事件：客户愿意把真实场景交给你，意味着产品开始走出实验室。

\begin{figure}[H]
\centering
\includegraphics[width=0.96\textwidth]{first-big-order.png}
\caption{第一笔 2000 万大单：大单让技术团队第一次被市场验证，也把交付压力放大。自制概念图，依据 00:49:02--00:55:45 对谈内容整理。}
\end{figure}

\begin{warningbox}{大单也会放大风险}
硬件大单如果交付失败，会同时伤害现金流、客户信任和行业口碑。签单只是开始，真正的考验是交付、稳定性和售后。
\end{warningbox}

\begin{knowledgebox}{硬科技早期验证链}
\begin{tabularx}{\textwidth}{p{0.20\textwidth}p{0.34\textwidth}X}
\toprule
阶段 & 验证什么 & 关键证据 \\
\midrule
样机 & 技术能不能跑起来 & 演示、测试数据、专家认可。 \\
小单 & 客户是否愿意试用 & 试点合同、现场反馈。 \\
大单 & 客户是否愿意承担风险 & 真实预算、真实交付压力。 \\
复购 & 产品是否真正有用 & 客户继续采购和扩规模。 \\
量产 & 能否进入产业链 & 稳定交付、质量体系、成本下降。 \\
\bottomrule
\end{tabularx}
\end{knowledgebox}

\begin{knowledgebox}{融资、订单、回款的区别}
\begin{tabularx}{\textwidth}{p{0.20\textwidth}p{0.34\textwidth}X}
\toprule
概念 & 证明什么 & 不能证明什么 \\
\midrule
融资 & 投资人相信未来可能性 & 不能证明客户一定买单。 \\
订单 & 客户愿意把预算给你 & 不能证明你能稳定交付。 \\
回款 & 客户接受交付并履约付款 & 不能证明规模化成本已解决。 \\
复购 & 客户愿意持续采购 & 才开始接近真实产品市场匹配。 \\
\bottomrule
\end{tabularx}
\end{knowledgebox}

\subsection{本章小结}

禾赛早期的组织和资本结构说明，硬科技创业要同时解决团队信任、融资节奏、客户订单和交付信用。技术只是进入牌桌的门票。

